% Part: turing-machines
% Chapter: machines-computations
% Section: well-behaved-machines

\documentclass[../../../include/open-logic-section]{subfiles}

\begin{document}

\olfileid{tur}{mac}{dis}
\olsection{پابند مشیناں}

\begin{explain}
\olref[hal]{sec} وچ
اسی اوہ ٹیورنگ مشیناں
ویکھیاں سن جنہاں دی
اکّو مقرر رُکّن والی
حالت~$h$ اے۔ جے ایہ
مشیناں کدی رُکّن،
تاں لازماً حالت~$h$
وچ ہی رُکّن گیاں۔
ایس طرح اکّو
رُکّن والی حالت
رکھن والیاں مشیناں
عام ٹیورنگ مشیناں
دے مقابلے وچ زیادہ
پابند نیں۔ مشیناں
دے عمل اُتے ہور
پابندیاں وی لا سکدے
آں۔ مثال دے طور
اُتے اسی مشین نوں
خانہ~$0$ اُتے پٹی
دے سِرے دا نشان
مٹاؤن توں، یا
خانہ~$0$ توں
کھبے پاسے جان دی
کوشش توں، منع
نہیں کیتا۔ (ساڈی
تعریف مطابق ایس
صورت وچ سر بس
خانہ~$0$ اُتے
ہی رہندا اے۔ کجھ ہور تعریفاں مطابق
ایس صورت وچ مشین رُک جاندی اے۔) ایہ وی کدے
کدے مفید ہوندا اے
کہ جے مشین رُکے
تاں خانہ~$1$
اُتے ہی رُکے۔
\end{explain}

\begin{defn}\ollabel{defn:disciplined}
ٹیورنگ مشین~$M$
\emph{پابند} اے اُدوں
تے صرف اُدوں جدوں:
\begin{enumerate}
    \item اوس دی اکّو
    مقرر رُکّن والی
    حالت~$h$ ہووے،
    \item جے اوہ رُکے
    تاں اوس دا سر خانہ~$1$
    پڑھ رہیا ہووے،
    \item اوہ $\TMendtape$
    نشان خانہ~$0$
    اُتے کدی نہ مٹاوے،
    تے
    \item اوہ خانہ~$0$
    توں کھبے پاسے
    جان دی کوشش کدی
    نہ کرے۔
\end{enumerate}
\end{defn}

\begin{explain}
اسی پہلے بیان کر
چکے آں کہ کسے وی
ٹیورنگ مشین نوں
اوہو عمل کرن والی اک
مشین وچ بدلیا
جا سکدا اے جیہدی
مقرر رُکّن والی حالت
ہووے۔ ایہ کم بس
نئی حالت~$h$
شامل کر کے، تے
ہر اوہو جیہے جوڑے لئی
$\delta(q, \sigma) = \tuple{h, \sigma, N}$
حکم شامل کر کے
ہوندا اے جتھے
$\tuple{q,\sigma}$ اُتے
اصل $\delta$~غیر
تعریف شدہ ہووے۔
ایہ سچ اے، اگرچہ
ثابت کرنا لمبا کم
اے، کہ ہر ٹیورنگ
مشین~$M$ نوں اک
پابند ٹیورنگ
مشین~$M'$ وچ
بدلیا جا سکدا اے
جیڑی اوہناں ہی
انپٹاں اُتے رُکے
تے اوہی آؤٹ پٹ
دیوے۔ مثال دے
طور اُتے، جے مشین
رُکّن لگے تے
خانہ~$1$ اُتے
نہ ہووے، تاں اسی
ہور حکم شامل کر
سکدے آں تاں جو
سر کھبے پاسے
چلدا رہے، جد تک
پٹی دے سِرے دا
نشان نہ لبھ لوے؛
فیر سر اک خانہ سجّے ول سرکے
تے مشین رُک جاوے۔ باقی شرطاں
کِویں پوری کیتیاں
جان، ایہ سوچݨ
لئی اسی تہانوں
چھڈدے آں۔
\end{explain}

\begin{ex}
    \olref{fig:adder-disc}
    وچ اسی
    \olref[una]{ex:adder}
    والی جمع دی مشین
    نوں پابند مشین
    وچ بدلدے آں۔
    \begin{figure}\[
\begin{tikzpicture}[->,>=stealth',shorten >=1pt,auto,node distance=2.8cm,
                    semithick]
  \tikzstyle{every state}=[fill=none,draw=black,text=black]

  \node[initial,state]         (A)              {$q_0$};
  \node[state]         (B) [right of=A] {$q_1$};
  \node[state]         (C) [below right of=B] {$q_2$};
  \node[state]         (D) [below left of=C] {$q_3$};
  \node[state]         (H) [left of=D] {$h$};

  \path (A) edge node {\TMtrans{\TMblank}{\TMstroke}{\TMstay}} (B)
           edge [loop below] node {\TMtrans{\TMstroke}{\TMstroke}{\TMright}} (B)
        (B) edge [loop below] node {\TMtrans{\TMstroke}{\TMstroke}{\TMright}} (B)
            edge node {\TMtrans{\TMblank}{\TMblank}{\TMleft}} (C)
        (C) edge node {\TMtrans{\TMstroke}{\TMblank}{\TMleft}} (D)
        (D) edge [loop above] node {\TMtrans{\TMstroke}{\TMstroke}{\TMleft}} (D)
        edge node {\TMtrans{\TMendtape}{\TMendtape}{\TMright}} (H);
\end{tikzpicture}
\]\caption{پابند جمع والی مشین}
\ollabel{fig:adder-disc}
\end{figure}
\end{ex}

\begin{prop}\ollabel{prop:disciplined}
ہر ٹیورنگ مشین~$M$
لئی اک پابند ٹیورنگ
مشین~$M'$ موجود اے
جیڑی آؤٹ پٹ~$O$
نال رُکدی اے جے
$M$~آؤٹ پٹ~$O$
نال رُکدی ہووے،
تے نہیں رُکدی
جے $M$~نہ رُکے۔
خاص طور اُتے،
جے فنکشن
$f\colon\Nat^n \to \Nat$
دا حساب کوئی
ٹیورنگ مشین
کردی اے، تاں
اوہدا حساب اک
پابند ٹیورنگ
مشین وی کردی
اے۔
\end{prop}

\begin{prob}\label{tur:mac:dis:prob:disc-succ}
فنکشن $f(x) = x+1$
دا حساب کرن والی
پابند مشین دیو۔
\end{prob}


\begin{prob}\label{tur:mac:dis:prob:copier}
ایسی پابند مشین
لبھو جیڑی
انپٹ $\TMstroke^n$
نال شروع ہو کے
آؤٹ پٹ
$\TMstroke^n \concat \TMblank \concat \TMstroke^n$
دیوے۔
\end{prob}

\end{document}
